% p103_2; source PDF page 103; manually verified editable table.
\begin{tikzpicture}[x=1cm,y=-1cm,inner sep=0pt,outer sep=0pt]
\path[use as bounding box] (0,0) rectangle (16.000,13.575);
\draw[line width=.4pt] (0.000,0.000) rectangle (3.100,0.855);
\node[align=center,text width=2.840cm,font={\linespread{1}\fontsize{10}{13.50}\selectfont\bfseries}] at (1.550,0.427) {反击对象};
\draw[line width=.4pt] (3.100,0.000) rectangle (9.700,0.855);
\node[align=center,text width=6.340cm,font={\linespread{1}\fontsize{10}{13.50}\selectfont\bfseries}] at (6.400,0.427) {结论和理由};
\draw[line width=.4pt] (9.700,0.000) rectangle (16.000,0.855);
\node[align=center,text width=6.040cm,font={\linespread{1}\fontsize{10}{13.50}\selectfont\bfseries}] at (12.850,0.427) {范例};
\draw[line width=.4pt] (0.000,0.855) rectangle (3.100,2.960);
\node[align=left,text width=2.840cm,font={\linespread{1}\fontsize{10}{13.50}\selectfont}] at (1.550,1.907) {正当防卫};
\draw[line width=.4pt] (3.100,0.855) rectangle (9.700,2.960);
\node[align=left,text width=6.340cm,font={\linespread{1}\fontsize{10}{13.50}\selectfont}] at (6.400,1.907) {对正当防卫的反击，不构成正当防卫，而是故意犯罪，因为正当防卫不属于不法侵害};
\draw[line width=.4pt] (9.700,0.855) rectangle (16.000,2.960);
\node[align=left,text width=6.040cm,font={\linespread{1}\fontsize{10}{13.50}\selectfont}] at (12.850,1.907) {甲杀乙，乙反击甲，甲又反击乙。甲的反击构成故意犯罪};
\draw[line width=.4pt] (0.000,2.960) rectangle (3.100,6.165);
\node[align=left,text width=2.840cm,font={\linespread{1}\fontsize{10}{13.50}\selectfont}] at (1.550,4.562) {紧急避险};
\draw[line width=.4pt] (3.100,2.960) rectangle (9.700,6.165);
\node[align=left,text width=6.340cm,font={\linespread{1}\fontsize{10}{13.50}\selectfont}] at (6.400,4.562) {（1）对紧急避险的反击，不构成正当防卫，因为紧急避险不属于不法侵害。（2）对紧急避险的反击不构成紧急避险，因为不符合法益衡量的要求。（3）对紧急避险的反击属于不法侵害};
\draw[line width=.4pt] (9.700,2.960) rectangle (16.000,6.165);
\node[align=left,text width=6.040cm,font={\linespread{1}\fontsize{10}{13.50}\selectfont}] at (12.850,4.562) {几条恶犬追咬甲，甲不得已闯入乙家。乙见状，将甲推出去。乙不构成紧急避险，因为乙保护的法益（住宅安宁状态）小于损害的法益（甲的身体健康）（2024 年试题）};
\draw[line width=.4pt] (0.000,6.165) rectangle (3.100,8.770);
\node[align=left,text width=2.840cm,font={\linespread{1}\fontsize{10}{13.50}\selectfont}] at (1.550,7.467) {假想防卫、事后防卫、防卫过当；假想避险、事后避险、避险过当};
\draw[line width=.4pt] (3.100,6.165) rectangle (9.700,8.770);
\node[align=left,text width=6.340cm,font={\linespread{1}\fontsize{10}{13.50}\selectfont}] at (6.400,7.467) {对这些行为的反击，构成正当防卫，因为这些行为都不构成正当防卫、紧急避险，都属于不法侵害};
\draw[line width=.4pt] (9.700,6.165) rectangle (16.000,8.770);
\node[align=left,text width=6.040cm,font={\linespread{1}\fontsize{10}{13.50}\selectfont}] at (12.850,7.467) {甲伤害乙后，离去，乙追上甲，用刀砍甲。甲反击乙。乙砍甲，属于事后防卫，属于不法侵害。甲反击乙构成正当防卫};
\draw[line width=.4pt] (0.000,8.770) rectangle (3.100,13.575);
\node[align=left,text width=2.840cm,font={\linespread{1}\fontsize{10}{13.50}\selectfont}] at (1.550,11.172) {偶然防卫、偶然避险};
\draw[line width=.4pt] (3.100,8.770) rectangle (9.700,13.575);
\node[align=left,text width=6.340cm,font={\linespread{1}\fontsize{10}{13.50}\selectfont}] at (6.400,11.172) {对其反击，观点展示：（1）按照防卫（避险）认识不要说，偶然防卫（偶然避险）构成正当防卫（紧急避险），不是不法侵害，因此对其不能进行正当防卫。（2）按照防卫（避险）认识必要说，偶然防卫（偶然避险）不构成正当防卫（紧急避险），客观上属于不法侵害，因此对其可以进行正当防卫};
\draw[line width=.4pt] (9.700,8.770) rectangle (16.000,13.575);
\node[align=left,text width=6.040cm,font={\linespread{1}\fontsize{10}{13.50}\selectfont}] at (12.850,11.172) {乙正要开枪杀丙，甲正要开枪杀乙，甲没有看到乙要杀人，甲只想杀死乙。甲的行为属于偶然防卫。丁反击甲。防卫认识不要说认为，丁不构成正当防卫。防卫认识必要说认为，丁构成正当防卫};
\end{tikzpicture}
